% Заполните значения, не меняя разметку титульного листа.
\newcommand{\ReportMinistry}{МИНОБРНАУКИ РОССИИ}
\newcommand{\ReportUniversity}{Федеральное государственное автономное образовательное учреждение высшего образования «Санкт-Петербургский государственный электротехнический университет «ЛЭТИ» им. В.~И.~Ульянова (Ленина)»}
\newcommand{\ReportDepartment}{[КАФЕДРА]}
\newcommand{\ReportDiscipline}{[ДИСЦИПЛИНА]}
\newcommand{\ReportTopic}{Решение задач линейного программирования с помощью симплекс-метода}
\newcommand{\ReportGroup}{[ГРУППА]}
\newcommand{\ReportStudentOne}{[ФИО СТУДЕНТА 1]}
\newcommand{\ReportStudentTwo}{[ФИО СТУДЕНТА 2]}
\newcommand{\ReportStudentThree}{[ФИО СТУДЕНТА 3]}
\newcommand{\ReportStudentFour}{[ФИО СТУДЕНТА 4]}
\newcommand{\ReportTeacher}{[ФИО ПРЕПОДАВАТЕЛЯ]}
\newcommand{\ReportYear}{[ГОД]}

\begin{titlepage}
  \thispagestyle{empty}
  \begin{center}
    \ReportMinistry

    \vspace{0.5\baselineskip}
    \ReportUniversity

    \vspace{0.5\baselineskip}
    Кафедра \ReportDepartment

    \vfill

    \textbf{ОТЧЕТ\\по практической работе № 1}

    \vspace{1.5\baselineskip}
    по дисциплине «\ReportDiscipline»

    \vspace{0.5\baselineskip}
    на тему «\ReportTopic»
  \end{center}

  \vfill

  \noindent
  \begin{tabularx}{\textwidth}{@{}X l l@{}}
    Студент группы \ReportGroup & \rule{3cm}{0.4pt} & \ReportStudentOne \\
    Студент группы \ReportGroup & \rule{3cm}{0.4pt} & \ReportStudentTwo \\
    Студент группы \ReportGroup & \rule{3cm}{0.4pt} & \ReportStudentThree \\
    Студент группы \ReportGroup & \rule{3cm}{0.4pt} & \ReportStudentFour \\
    & {\small подписи} & \\
    Проверил & \rule{3cm}{0.4pt} & \ReportTeacher \\
    & {\small подпись} &
  \end{tabularx}

  \vfill

  \begin{center}
    Санкт-Петербург\\
    \ReportYear
  \end{center}
\end{titlepage}
